% Einheit 3 von 4 – Senkrecht zu Geraden und Ebenen (bank/orthogonalitaet/e3.jsonl, zusammenbau v0.1)
%% TODO Zeitmarke Einheit 3: Marken-Zeile nicht lesbar
%% TODO Zweigzeile Teil 1: was hier gelernt wird (Satz in Schülersprache aus dem Einheitstitel) – Einheit 3
\einheitenkopf[e3]{Einheit 3 von 4 · Senkrecht zu Geraden und Ebenen}
\zweigzeile{Abitur GK}
\setcounter{aufgabe}{15}

%% TODO Titel: Ich-kann-Satz fehlt in der Bank (Platzhalter: Kettenname) – Nr. 16
%% TODO Anweisung: ein Satz über der ersten Teilaufgabe fehlt in der Bank – Nr. 16
\begin{aufgabe}{Senkrecht zu Geraden und Ebenen}
%% TODO \rechenplatz{n}: Schrittzahl des Grundfalls fehlt in der Bank (nur bei fester Schreibform, 2.3 b)
\begin{teile}
\teil Aufgabe: „Zeige, dass die Gerade $h$ senkrecht auf der Ebene $E$ steht.“ Welche Prüfregel? \\ \kreuz{zwei Richtungen – Skalarprodukt null} \\ \kreuz{Gerade gegen Ebene – Richtungsvektor kollinear zum Normalenvektor} \\ \kreuz{Ebene gegen Ebene – Skalarprodukt der Normalenvektoren null}
\teil Gegeben sind $g\colon \vec x = (2 | 1 | 5) + r \cdot (2 | -1 | 3)$ und $E\colon 2x - y + 3z = 4$. Begründe, dass $g$ senkrecht zu $E$ steht \hfill (Abitur 2022 LK).
\teil Ein Prisma $ABCDEF$ hat die Grundfläche $ABC$ in der Ebene $L\colon 2x - y + 2z = 6$ mit $A(3 | 0 | 0)$, $B(1 | 2 | 3)$, $C(0 | -2 | 2)$; die Seitenkante $BE$ führt zu $E(5 | 0 | 7)$. Begründe, dass das Prisma gerade ist \hfill (Abitur 2020 GK).
\teil Gegeben sind $E\colon \vec x = (2 | 0 | 1) + r \cdot (2 | 1 | 2) + s \cdot (1 | 0 | -1)$ und der Vektor $\vec n = (1 | -4 | 1)$. Weise nach, dass $\vec n$ senkrecht zu $E$ steht \hfill (Abitur 2025 GK).
\teil Gegeben ist $g\colon \vec x = (5 | 2 | -1) + s \cdot (-3 | 0 | 2)$; die Gerade $h$ verläuft parallel zur $y$-Achse und schneidet $g$ im Punkt $(5 | 2 | -1)$. Untersuche, ob $g$ und $h$ senkrecht zueinander verlaufen \hfill (Abitur 2025 GK).
\teil Gegeben sind $E\colon 4x - 3y + z = 2$ und $F\colon x + 2y + kz = 5$. Bestimme $k$ so, dass $F$ senkrecht zu $E$ steht \hfill (Abitur 2022 GK).
\teil Gegeben sind $E\colon x + 2y - z = 3$ und $F\colon 3x - y + 2z = 1$; ihre Schnittgerade hat den Richtungsvektor $(3 | -5 | -7)$. Gib eine Gleichung einer Ebene $H$ an, die senkrecht zu $E$ und zu $F$ steht \hfill (Abitur 2021 GK).
\end{teile}
\end{aufgabe}

%% TODO Titel: Ich-kann-Satz fehlt in der Bank (Platzhalter: Kettenname) – Nr. 17
%% TODO Anweisung: ein Satz über der ersten Teilaufgabe fehlt in der Bank – Nr. 17
\begin{aufgabe}{Senkrecht zu Geraden und Ebenen – Fehler finden, Begründen, Anwendung}
\begin{teile}
\teil Für $g\colon \vec x = (2 | 1 | 5) + r \cdot (2 | -1 | 3)$ und $E\colon 2x - y + 3z = 4$ rechnet Paul $(2 | -1 | 3) \circ (2 | -1 | 3) = 14 \ne 0$ und folgert, $g$ stehe nicht senkrecht auf $E$. Finde den Fehler.
\teil Begründe, warum eine Gerade genau dann senkrecht auf einer Ebene steht, wenn ihr Richtungsvektor ein Vielfaches des Normalenvektors ist.
\teil Ein Hang wird durch die Ebene $E\colon x - 4z = -8$ beschrieben; ein Mast steht von $F(0 | 3 | 2)$ aus in Richtung $(2 | 0 | -8)$ (1 LE = 1 m). Steht der Mast senkrecht auf dem Hang? Begründe.
\end{teile}
\end{aufgabe}
